\subsection{A fixed finite trace diagram cannot contain every power}\label{sec:fixed-diagram-obstruction}

The following elementary obstruction applies whenever a proposed realization uses a fixed finite family of vector spaces and endomorphisms. It permits arbitrary insertions, including ones that do not commute with those endomorphisms.

\begin{proposition}\label{prop:fixed-diagram-obstruction}
Fix $q>1$, finitely many finite-dimensional complex vector spaces $V_\ell$, endomorphisms $F_\ell$, and scalars $\mu_\ell$. Fix $c\ne0$ and $\eta\in\mathbb C$. There do not exist endomorphisms $A_{k,\ell}$, chosen separately for each $k\ge1$, such that
\[
 \sum_\ell\mu_\ell\operatorname{Tr}(A_{k,\ell}F_\ell^n)
 =\eta+c(q^n+2+q^{-n})^k
 \qquad(n,k\ge1).
\]
In particular the obstruction is unchanged by adding a fixed weight-zero line to every right side.
\end{proposition}

\begin{proof}
For every fixed $k$, the generating function of the left side is
\[
 \sum_\ell\mu_\ell\operatorname{Tr}
 \bigl(A_{k,\ell}zF_\ell(I-zF_\ell)^{-1}\bigr).
\]
It is rational, with finite poles among the inverses of the nonzero eigenvalues of the fixed matrices $F_\ell$. This is a finite set independent of $k$. No commutation hypothesis is used in the matrix geometric series.

The binomial identity
\[
 (t+2+t^{-1})^k=\sum_{j=-k}^k\binom{2k}{k+j}t^j
\]
shows that the proposed right-side generating function is
\[
 \eta\frac{z}{1-z}
 +c\sum_{j=-k}^k\binom{2k}{k+j}
       \frac{q^jz}{1-q^jz}.
\]
It has an actual pole at $z=q^{-k}$: only its $j=k$ summand is singular there, with nonzero coefficient $c$. The fixed weight-zero term is analytic there. Equality of all coefficients would imply equality of the rational functions, so the fixed finite pole set would have to contain the distinct numbers $q^{-k}$ for every $k\ge1$. This is impossible.
\end{proof}

The proof uses fixed pole locations, and does not assert that inserted traces have only simple poles. Plain traces satisfy $\operatorname{Tr}(F^n)=\sum_\lambda m_\lambda\lambda^n$, because the nilpotent parts of Jordan blocks have zero trace. In contrast,
\[
 F=\begin{pmatrix}\lambda&1\\0&\lambda\end{pmatrix},\qquad
 A=\begin{pmatrix}0&0\\1&0\end{pmatrix},\qquad\lambda\ne0,
\]
give $\operatorname{Tr}(AF^n)=n\lambda^{n-1}$ and the double-pole generating function $z/(1-\lambda z)^2$.

A finite support diagram of finite-dimensional complex vector spaces, equipped with compatible Frobenius operators, gives a finite family of this kind. A M\"obius trace on its support poset uses fixed coefficients $\mu_\ell$. Thus neither such a fixed diagram nor the addition of a rational basepoint can realize all the displayed powers simultaneously. Separate realizations whose spaces depend on the word remain possible. The diagram interpretation uses the inverse additive structure of semimodules in \citet[Proposition~6.1]{FGJM}; Frobenius operators are additional data, not consequences of zero adjunction.
